\documentclass[9pt,twoside]{extarticle}
\usepackage[a4paper,top=14mm,bottom=14mm,left=16mm,right=16mm,headheight=12pt,headsep=3.5mm,footskip=7mm,includeheadfoot]{geometry}
\usepackage{fontspec}
\defaultfontfeatures{Ligatures=TeX,Scale=MatchLowercase}
\setmainfont{Libertinus Serif}
\setsansfont{Libertinus Sans}
\usepackage{mathtools}
\usepackage{array}
\usepackage{unicode-math}
\setmathfont{Libertinus Math}
\usepackage[english]{babel}
\usepackage{microtype}
\usepackage{adjustbox}
\usepackage{fancyhdr}
\usepackage{ragged2e}
\usepackage[hidelinks,unicode]{hyperref}
\usepackage{bookmark}
\hypersetup{pdftitle={Local Constants in the Functional Equation of the Artin L-Function of an Orthogonal Representation},pdfauthor={Pierre Deligne},pdfsubject={Standalone faithful English edition},pdfkeywords={Deligne, Artin L-function, local constants, orthogonal representations, editable LaTeX}}
\pagestyle{fancy}
\fancyhf{}
\fancyhead[LE,RO]{\footnotesize\sffamily Local constants - English edition}
\fancyhead[LO,RE]{\footnotesize\sffamily Printed folio \thepage}
\fancyfoot[C]{\footnotesize\thepage}
\renewcommand{\headrulewidth}{0.25pt}
\renewcommand{\footrulewidth}{0pt}
\setlength{\parindent}{1em}
\setlength{\parskip}{0.13em}
\setlength{\abovedisplayskip}{3.2pt plus 1pt minus 1pt}
\setlength{\belowdisplayskip}{3.2pt plus 1pt minus 1pt}
\setlength{\abovedisplayshortskip}{2pt}
\setlength{\belowdisplayshortskip}{2pt}
\setlength{\jot}{1.5pt}
\emergencystretch=1.4em
\providecommand{\dashrightarrow}{\mathrel{⇢}}
\newcommand{\EditionHeading}[1]{\par\medskip{\noindent\centering\large\bfseries #1\par}\smallskip}
\newcommand{\EditionSection}[1]{\par\medskip{\noindent\large\bfseries #1\par}\smallskip}
\newcommand{\EditionPage}[3]{%
  \ifnum#1>1\clearpage\fi
  \pdfbookmark[1]{Printed folio #2}{en-folio-#2}%
  \hypertarget{en-page-#1}{}%
  \noindent\begin{adjustbox}{max totalsize={\textwidth}{0.915\textheight},center}
  \begin{minipage}[t]{\textwidth}
  \fontsize{8.45}{10.05}\selectfont\justifying
  #3
  \end{minipage}
  \end{adjustbox}%
}
\begin{document}
\setcounter{page}{299}
\pdfbookmark[0]{Local Constants in the Functional Equation of the Artin L-Function of an Orthogonal Representation}{d026-title-en}

\EditionPage{1}{299}{%
\begin{center}\fontsize{15}{17}\selectfont\bfseries Local Constants in the Functional Equation\\of the Artin L-Function of an Orthogonal Representation\end{center}

\begin{center}\large Pierre Deligne (Bures-sur-Yvette)\end{center}

\begin{center}\itshape To Jean-Pierre Serre, in token of admiration.\end{center}

\EditionHeading{Contents}

{\footnotesize\setlength{\parskip}{0pt}1. Statement of the theorem ................................... 299\par
2. Lemmas on virtual representations of finite groups .......... 302\par
3. Proof of 1.5 ................................................ 307\par
4. Weil groups and the Hasse reciprocity law .................... 309\par
5. Orthogonal representations of Weil groups .................... 313\par}

\EditionSection{1. Statement of the theorem}

(1.1) Let \(K\) be a global field (an A-field in Weil’s sense), let \(K'\) be a finite Galois extension of \(K\), and let \(V\) be a complex representation of the Galois group \(\operatorname{Gal}(K'/K)\). We follow the conventions of [1] in defining the corresponding Artin \(L\)-function: the factors at infinity are included, and the local nonarchimedean factors are defined in terms of geometric Frobenius elements (the inverses of Frobenius substitutions). We further normalize the reciprocity isomorphism of local class field theory so that it sends geometric Frobenius elements to uniformizers (cf. [1] 3.6). These sign questions are in fact immaterial: our essential calculations will be made modulo 2. The function \(L(V,s)\) satisfies a functional equation\par

\[
\tag{1.1.1} L(V,s)=\varepsilon(V,s)L(V^*,1-s),
\]

where \(\varepsilon(V,s)\) is the product of a constant and an exponential factor, and \(V^*\) is the dual of \(V\).\par

Let \(f\) be the norm of the Artin conductor of \(V\), and let \(D\) be the absolute value of the discriminant of \(K\) (respectively \(q^{2g-2}\) for a function field); define\par

\[
\tag{1.1.2} W(V)=\varepsilon\!\left(V,\tfrac12\right).
\]

Then\par

\[
\tag{1.1.3} \varepsilon(V,s)=W(V)\bigl(f\cdot D^{\dim(V)}\bigr)^{\frac12-s},
\]

\[
\tag{1.1.4} W(V)\cdot W(V^*)=1.
\]
}

\EditionPage{2}{300}{%
If \(V\) is isomorphic to \(V^*\), i.e. if the character of \(V\) is real-valued, (1.1.4) shows that \(W(V)=\pm1\). Fröhlich and Queyrut [2] showed that if \(V\) is in fact orthogonal, i.e. realizable over \(\mathbf R\), then \(W(V)=+1\).\par

(1.2) Let \(K\) be a local field (= a nondiscrete locally compact field), let \(K'\) be a finite Galois extension of \(K\) contained in a separable closure \(\bar K\) of \(K\), let \(V\) be a complex representation of the Galois group \(G=\operatorname{Gal}(K'/K)\), let \(\psi:K\to\mathbf C^*\) be a nontrivial additive character, and let \(dx\) be a Haar measure on \(K\). In fact, \(K'\) plays no role: all that is used is the representation of \(\operatorname{Gal}(\bar K/K)\) obtained from \(V\) by inflation; it is again denoted by \(V\). I refer to [1] §§ 4, 5 for the definition of \(\varepsilon(V,\psi,dx,s)\). In [1], more generally, \(V\) is taken to be a representation of the Weil group \(W(\bar K/K)\) (of which \(\operatorname{Gal}(\bar K/K)\) is a completion in the nonarchimedean case and a quotient in the archimedean case); one defines \(\varepsilon(V,\psi,dx)\), and sets \(\varepsilon(V,\psi,dx,s)=\varepsilon(V\otimes\omega_s,\psi,dx)\), where \(\omega_s\) is the quasicharacter \(\lVert x\rVert^s\) of \(K^*\simeq W(\bar K/K)^{\mathrm{ab}}\). As a function of \(s\), \(\varepsilon(V,\psi,dx,s)\) is the product of a constant and an exponential factor.\par

Starting from a global situation as in 1.1, let \(K_v\) run through the completions of \(K\), let \(V_v\) be the restriction of \(V\) to a decomposition group at \(v\), let the \(\psi_v\) be the local components of a character of \(A_K/K\), and suppose that\par

\[
\int_{A_K/K}\bigotimes_v dx_v=1.
\]

Then\par

\[
\tag{1.2.1} \varepsilon(V,s)=\prod_v\varepsilon(V_v,\psi_v,dx_v,s).
\]

Both locally and globally, \(\varepsilon\) is multiplicative in \(V\): \(\varepsilon(V\oplus V')=\varepsilon(V)\cdot\varepsilon(V')\). This makes it possible to define \(\varepsilon\) for virtual representations (elements of the Grothendieck group \(R(G,\mathbf C)\) of the category of representations of \(G\)).\par

Return to the local case. If \(V\) is a virtual representation of dimension 0, \(\varepsilon(V,\psi,dx,s)\) is independent of \(dx\), which is omitted from the notation ([1] 5.3). If \(V\) has dimension 0 and determinant \(\chi:K^*\simeq W(\bar K/K)^{\mathrm{ab}}\to\mathbf C^*\), then ([1] 5.4)\par

\[
\tag{1.2.2} \varepsilon(V,\psi(ax),s)=\chi(a)\varepsilon(V,\psi,s).
\]

In particular, if \(\chi=1\), \(\varepsilon(V,\psi,s)\) is independent of \(\psi\); \(\psi\) is then omitted from the notation.\par

The Artin conductor \(f\) of \(V\) is by definition 1 when \(K\) is archimedean. When \(K\) is nonarchimedean, it is the number denoted \(q^{a(V)}\) in [1] (\(q=\) the number of elements of the residue field). For \(V\) of dimension 0 and determinant 1, set\par

\[
\tag{1.2.3} W(V)=\varepsilon\!\left(V,\tfrac12\right).
\]

Then\par

\[
\tag{1.2.4} \varepsilon(V,s)=W(V)f^{\frac12-s},
\]

\[
\tag{1.2.5} W(V)\cdot W(V^*)=1.
\]

If \(V\) is isomorphic to its dual, i.e. if its character is real, (1.2.5) shows that \(W(V)=\pm1\).\par

The group \(R(G,\mathbf R)\) of real virtual representations of \(G\), the Grothendieck group of the category of real representations of \(G\), is a subgroup of the group of virtual representations of \(G\) having real character. The elements of\par
}

\EditionPage{3}{301}{%
this subgroup are said to be realizable over \(\mathbf R\). We propose to calculate \(W(V_v)\) when \(V_v\) has dimension 0, determinant 1, and is realizable over \(\mathbf R\).\par

(1.3) Let \(G\) be a finite group. A real representation \(\rho:G\to GL(V)\) has a Stiefel-Whitney class\par

\[
w^*(V)=1+w^1(V)+\cdots\in H^*(G,\mathbf Z/2)=\prod_i H^i(G,\mathbf Z/2).
\]

It may be defined as follows: the representation \(V\) defines a vector bundle (indeed, a local system) \(\mathcal V\) on the classifying space \(BG\) of the group \(G\), and \(w^*(V)\) is the Stiefel-Whitney class of \(\mathcal V\) in \(H^*(BG,\mathbf Z/2)=H^*(G,\mathbf Z/2)\). The same construction applies to real virtual representations: \(w^*\) is the composite homomorphism\par

\[
R(G,\mathbf R)\longrightarrow KO(BG)\overset{w^*}{\longrightarrow}H^*(BG,\mathbf Z/2)^\times=H^*(G,\mathbf Z/2)^\times.
\]

Via the isomorphism \(H^1(G,\mathbf Z/2)=\operatorname{Hom}(G,\{\pm1\})\), \(w^1(V)\) identifies with the character \(g\mapsto\det\rho(g)\). The group \(H^2(G,\mathbf Z/2)\) classifies central extensions of \(G\) by the two-element group. If \(w^1(V)=0\), then \(w^2(V)\) is the class of the extension obtained by pulling back under \(\rho\) the double cover \(\operatorname{Spin}(V,Q)\) of \(SO(V,Q)\), where \(Q\) is any positive-definite \(G\)-invariant quadratic form on \(V\).\par

(1.4) Retain the notation of 1.2, and let \(\operatorname{cl}(w^2(V))\) be the image of the class \(w^2(V)\in H^2(G,\mathbf Z/2)\) under the composite map\par

\[
\tag{1.4.1}
H^2(G,\mathbf Z/2)\overset{\mathrm{infl}}{\longrightarrow}H^2(\operatorname{Gal}(\bar K/K),\mathbf Z/2)
\overset{(1)}{\longrightarrow}H^2(\operatorname{Gal}(\bar K/K),\bar K^*)
\overset{\mathrm{inv}}{\longrightarrow}\mathbf Q/\mathbf Z,
\]

where (1) is induced by \(n\mapsto(-1)^n:\mathbf Z/2\to\bar K^*\), and where \(\mathrm{inv}\) is supplied by local class field theory. If \(K\) is complex or has characteristic 2, then \(\operatorname{cl}w^2(V)=0\). In the other cases, (1.4.1) identifies \(H^2(\operatorname{Gal}(\bar K/K),\mathbf Z/2)\) with \(\{0,\tfrac12\}\subset\mathbf Q/\mathbf Z\).\par

Our principal result is the following:\par

(1.5) Theorem. If \(V\) is a real virtual representation of dimension 0 and determinant 1, then\par

\[
W(V)=\exp\bigl(2\pi i\,\operatorname{cl}(w^2(V))\bigr).
\]

Thus, in characteristic 2, \(W(V)=1\). In characteristic \(\ne2\), \(W(V)=1\) or \(-1\) according as \(w^2(V)\) is or is not trivial in \(H^2(\operatorname{Gal}(\bar K/K),\mathbf Z/2)\).\par

(1.6) From this theorem one can deduce the theorem of Queyrut and Fröhlich:\par

(a) For \(V\) a real virtual representation of dimension 0 and determinant 1 over a global field \(K\) of characteristic \(\ne2\), the \(\operatorname{cl}(w^2(V_v))\) are the local invariants of \(w^2(V)\in H^2(\operatorname{Gal}(\bar K/K),\mathbf Z/2)\), and the reciprocity law ensures that their sum in \(\mathbf Q/\mathbf Z\) is zero; hence the product of the \(W(V_v)\) is one. In equal characteristic 2, all the \(W(V_v)\) are 1.\par

(b) If \(V\) has dimension one and is defined by a character of order two (respectively, is trivial), the corresponding \(L\)-function is a quotient of two zeta functions (respectively, a zeta function), and one uses the fact that the global constant of a zeta function is \(+1\).\par

(c) The conclusion follows by observing that every real virtual representation is a sum of representations of types (a) and (b).\par
}

\EditionPage{4}{302}{%
(1.7) We shall also prove the variants of these results obtained by replacing Galois groups by Weil groups (§ 5).\par

\EditionSection{2. Lemmas on real representations of finite groups}

(2.1) Proposition. Let \(G\) be a finite group, \(H\) a subgroup, \(V\) a real virtual representation of \(H\), and \(\operatorname{Tr}:H^2(H,\mathbf Z/2)\to H^2(G,\mathbf Z/2)\) the trace morphism (corestriction). If the dimension of \(V\) is divisible by 4 and its determinant is 1 (\(w^1(V)=0\)), then\par

\[
w^1\!\left(\operatorname{Ind}_H^G(V)\right)=0
\quad\text{and}\quad
w^2\!\left(\operatorname{Ind}_H^G(V)\right)=\operatorname{Tr}w^2(V).
\]

Let us translate this statement into topological terms. Let \(B_G\) be a classifying space for \(G\), let \(E_G\) be the universal cover of \(B_G\) (the universal principal homogeneous \(G\)-space), and set \(B_H=E_G/H\). This is a classifying space for \(H\). Let \(f\) denote the projection \(f:B_H\to B_G\). The diagrams\par

\[
\begin{array}{ccc}
R(H,\mathbf R)&\longrightarrow&KO(B_H)\\
\big\downarrow\scriptstyle{\operatorname{Ind}}&&\big\downarrow\scriptstyle{f_*}\\
R(G,\mathbf R)&\longrightarrow&KO(B_G)
\end{array}
\qquad\text{and}\qquad
\begin{array}{ccc}
H^*(B_H,\mathbf Z/2)&=&H^*(H,\mathbf Z/2)\\
\big\downarrow\scriptstyle{\operatorname{tr}}&&\big\downarrow\scriptstyle{f_*}\\
H^*(B_G,\mathbf Z/2)&=&H^*(G,\mathbf Z/2)
\end{array}
\]

commute (the definition of the maps \(f_*\) is recalled below), so that 2.1 follows from the following proposition.\par

(2.2) Proposition. Let \(f:X\to Y\) be a finite covering of topological spaces. If \(v\in KO(X)\) has dimension divisible by 4 and zero first Stiefel-Whitney class, then \(w^1(f_*v)=0\) and\par

\[
w^2(f_*v)=f_*(w^2(v)).
\]

Recall the definition of the maps \(f_*\).\par

(a) In K-theory: \(f_*\) is induced by the operation that associates to a vector bundle \(\mathcal V\) on \(X\) the vector bundle \(f_*\mathcal V\) on \(Y\) whose fiber at \(y\in Y\) is the direct sum\par

\[
\bigoplus_{f(x)=y}\mathcal V_x.
\]

(b) In cohomology. \(f_*\) is the composite of the isomorphism \(H^*(X,\mathbf Z/2)=H^*(Y,f_*\mathbf Z/2)\) and the morphism induced functorially by the trace morphism:\par

\[
\operatorname{Tr}_f:f_*\mathbf Z/2_X\longrightarrow\mathbf Z/2_Y:\ \text{“sum of the values over the sheets.”}
\]

We reduce to the case in which \(f\) has a constant number \(k\) of sheets and \(v\) is defined by an orthogonal vector bundle \(\mathcal V\) of constant dimension \(n\). Let \(P'\) be the corresponding \(O(n)\)-torsor and let \([P']\) be its class in \(H^1(X,O(n))\). The descriptions in 1.3 of \(w^1\) and \(w^2\) have the following analogue.\par

(a) \(w^1\in H^1(X,\mathbf Z/2)\) is the image of \([P']\) under\par

\[
\det:H^1(X,O(n))\longrightarrow H^1(X,\{\pm1\})=H^1(X,\mathbf Z/2).
\]

If \(w^1=0\), \(P'\) lifts to an \(SO(n)\)-torsor \(P\), of class \([P]\).\par
}

\EditionPage{5}{303}{%
(b) The exact sequence\par

\[
0\longrightarrow\mathbf Z/2\longrightarrow\operatorname{Spin}(n)\longrightarrow SO(n)\longrightarrow0
\]

induces \(\partial:H^1(X,SO(n))\to H^2(X,\mathbf Z/2)\), and, if \(w^1=0\), \(w^2=\partial[P]\).\par

The construction \(\mathcal V\mapsto f_*\mathcal V\) associates to an \(O(n)\)-torsor on \(X\) an \(O(nk)\)-torsor on \(Y\). It factors as follows: the local isomorphisms from the constant orthogonal vector bundle \((\mathbf R^n)^k\) to \(f_*\mathcal V\) which, at each point \(y\), carry each of the \(k\) factors of \((\mathbf R^n)^k\) to one of the factors in the decomposition\par

\[
(f_*\mathcal V)_y=\bigoplus_{f(x)=y}\mathcal V_x,
\]

identify with the local sections of a torsor under the semidirect product \(O(n)^k\rtimes\mathfrak S_k\) of \(O(n)^k\) by the symmetric group on \(k\) letters. The required \(O(nk)\)-torsor is obtained from this torsor and the morphism \(O(n)^k\rtimes\mathfrak S_k\to O(nk)\).\par

If \(w^1(\mathcal V)=0\), the structure group of \(f_*\mathcal V\) can be reduced to \(SO(n)^k\rtimes\mathfrak S_k\). If \(n\) is even, this group is contained in \(SO(nk)\), and \(w^1(f_*\mathcal V)=0\).\par

Suppose that the structure group of \(\mathcal V\) has been reduced to \(SO(n)\), giving an \(SO(n)\)-torsor \(P\) and an \(SO(n)^k\rtimes\mathfrak S_k\)-torsor \(Q\) on \(Y\), of class \([Q]\). Consider the exact sequence\par

\[
0\longrightarrow(\mathbf Z/2)^k\longrightarrow\operatorname{Spin}(n)^k\rtimes\mathfrak S_k\longrightarrow SO(n)^k\rtimes\mathfrak S_k\longrightarrow0.
\]

The quotient \(SO(n)^k\rtimes\mathfrak S_k\) acts on \((\mathbf Z/2)^k\) through \(\mathfrak S_k\), and the local system \(((\mathbf Z/2)^k)^Q\) on \(Y\) identifies with \(f_*\mathbf Z/2\). Via the isomorphism \(H^*(X,\mathbf Z/2)\simeq H^*(Y,f_*\mathbf Z/2)\), the class \(w^2(\mathcal V)=\partial[P]\) identifies with \(\partial[Q]\in H^2(Y,((\mathbf Z/2)^k)^Q)\), as is shown by a cocycle computation.\par

The proposition now follows from the existence, when \(n\) is divisible by 4, of a middle arrow making the diagram\par

\[
\tag{2.2.1}
\begin{array}{ccccccccc}
0&\to&(\mathbf Z/2)^k&\to&\operatorname{Spin}(n)^k\rtimes\mathfrak S_k&\to&SO(n)^k\rtimes\mathfrak S_k&\to&0\\
&&\big\downarrow\scriptstyle{\Sigma}&&\big\downarrow&&\big\downarrow\\
0&\to&\mathbf Z/2&\to&\operatorname{Spin}(nk)&\to&SO(nk)&\to&0
\end{array}
\]

commute, where \(\Sigma\) is the sum-of-coordinates homomorphism.\par

To construct (2.2.1), it is enough to construct a lift\par

\[
\begin{array}{ccc}
\mathfrak S_k&\longrightarrow&SO(nk)\\
&\searrow\ \dashrightarrow&\big\uparrow\\
&&\operatorname{Spin}(nk).
\end{array}
\]

The obstruction to constructing it is defined for even \(n\), and lies in \(H^2(\mathfrak S_k,\mathbf Z/2)\). It is additive in \(n\), because the diagram\par

\[
\begin{array}{ccccccccc}
0&\to&\mathbf Z/2\times\mathbf Z/2&\to&\operatorname{Spin}(n'k)\times\operatorname{Spin}(n''k)&\to&SO(n'k)\times SO(n''k)&\to&0\\
&&\big\downarrow\scriptstyle{a+b}&&\big\downarrow&&\big\downarrow\\
0&\to&\mathbf Z/2&\to&\operatorname{Spin}((n'+n'')k)&\to&SO((n'+n'')k)&\to&0
\end{array}
\]

commutes. It is therefore zero when \(n\) is divisible by 4.\par
}

\EditionPage{6}{304}{%
(2.3) Lemma. The real virtual representations of dimension 0 and determinant 1 of a finite group \(G\) form an ideal of \(R(G,\mathbf R)\). This ideal is stable under inflation, restriction, and induction.\par

That they form an ideal follows from the fact that \(\dim:R(G,\mathbf R)\to\mathbf Z\) is a homomorphism, that \(\det(V+W)=\det(V)\cdot\det(W)\), and that\par

\[
\det(V\otimes W)=\det(V)^{\dim(W)}\cdot\det(W)^{\dim(V)}.
\]

Stability under inflation and restriction is immediate, while stability under induction follows from 2.2 (the first assertion).\par

An analogous lemma holds for complex representations. It is proved in the same way ([1] 1.2).\par

(2.4) Let \(D_n\) be the dihedral group, the semidirect product of \(\mathbf Z/2\) by \(\mathbf Z/n\), and let \(\chi:\mathbf Z/n\to\mathbf C^*\) be a character. The representation \(\operatorname{Ind}(\chi)\) is realizable over \(\mathbf R\): it is realized on the real vector space \(\mathbf C\), on which \(\mathbf Z/n\) acts by \(\chi\) and \(\mathbf Z/2\) acts by complex conjugation. Its determinant is independent of \(\chi\). The virtual representation\par

\[
r(\chi)=\operatorname{Ind}_{\mathbf Z/n}^{D_n}([\chi]-[1])
\]

is therefore realizable over \(\mathbf R\), has dimension 0, and has trivial determinant.\par

(2.5) Proposition. Let \(G\) be a finite group. Every real virtual representation of \(G\) of dimension 0 and trivial determinant is a sum\par

(a) of the realification of a complex virtual representation of dimension 0, and\par

(b) of a \(\mathbf Z\)-linear combination of representations induced from representations \(r(\chi)\) of dihedral quotients of subgroups of \(G\).\par

Recall that the “realification” of a complex representation \(V\) of a group \(G\) is the representation of \(G\) on the same space \(V\), regarded as a real vector space by restriction of scalars from \(\mathbf C\) to \(\mathbf R\).\par

(2.6) Lemma. Let \(G\) be a finite group having a nilpotent subgroup \(G'\) of index 1 or 2. Every real virtual representation of \(G\) is a sum\par

(a) of the realification of a complex virtual representation of \(G\),\par

(b) of a \(\mathbf Z\)-linear combination of representations induced from representations \(r(\chi)\),\par

(c) of a \(\mathbf Z\)-linear combination of one-dimensional real representations.\par

Let us deduce the proposition from the lemma.\par

Brauer’s induction theorem (as generalized by Witt, J. Crelle, 190 (1952); see, for example, Curtis-Reiner, Representation theory of finite groups and associative algebras, Intersc. Publ., N.Y. 1962) asserts that the representation 1 of \(G\) is a \(\mathbf Z\)-linear combination \(\sum_i n_i\operatorname{Ind}_{H_i}^G(\rho_i)\) of representations induced from real representations of subgroups \(H_i\) of \(G\), each \(H_i\) being a semidirect product of a \(p\)-group \(P\) by a cyclic group \(C\), with \(P\) acting on \(C\) only by \(\pm1\). Lemma 2.3 and the identity\par

\[
V=V\otimes1=V\otimes\sum_i n_i\operatorname{Ind}_{H_i}^G(\rho_i)=\sum_i n_i\operatorname{Ind}_{H_i}^G(V\otimes\rho_i)
\]
}

\EditionPage{7}{305}{%
bring us back to the case where \(G\) itself is such a semidirect product, hence of the type considered in 2.6. For a virtual \(V\) of dimension 0 and determinant 1, consider the decomposition 2.6\par

\[
V=A+B+C=(A-\dim A\cdot 1)+B+(C+\dim A\cdot 1).
\]

The representation \((A-\dim A\cdot 1)\) is of type (a) in the proposition, and \(B\) is of type (b). Since \(V\), \((A-\dim A\cdot 1)\), and \(B\) have dimension 0 and determinant 1, \((C+\dim A\cdot 1)\) has the same properties; this reduces us to proving the proposition for \((C+\dim A\cdot 1)\), hence for a \(\mathbf Z\)-linear combination of one-dimensional real representations. Thus let \(V\) be of the form\par

\[
V=\sum_{i=1}^{k}[\varepsilon_i]-\sum_{i=1}^{l}[\eta_i]
\]

with \(\varepsilon_i\) and \(\eta_i\) characters of order dividing 2, \(k=l\) (dimension 0), and \(\prod\varepsilon_i=\prod\eta_i\) (determinant 1).\par

We have\par

\[
\begin{aligned}
\sum_{i=1}^{k}[\varepsilon_i]
&=k\cdot 1+\sum_{i=1}^{k}([\varepsilon_i]-1)\\
&=k\cdot 1-\sum_{i=1}^{k}([\varepsilon_i]-1)\otimes\left(\left[\prod_{j<i}\varepsilon_j\right]-1\right)
 +\sum_{i=1}^{k}([\varepsilon_i]-1)\otimes\left[\prod_{j<i}\varepsilon_j\right]\\
&=k\cdot 1+\left(\left[\prod\varepsilon_i\right]-1\right)
 -\sum_{i=1}^{k}([\varepsilon_i]-1)\left(\left[\prod_{j<i}\varepsilon_j\right]-1\right).
\end{aligned}
\]

It follows that \(V\) has an expression of the form\par

\[
V=\sum\pm\bigl([\alpha_i\beta_i]-[\alpha_i]-[\beta_i]+1\bigr),
\]

where the \(\alpha_i\) and \(\beta_i\) have order dividing 2. The terms in this sum are zero if \(\alpha_i=1\) or \(\beta_i=1\), are of type (a) if \(\alpha_i=\beta_i\), and are of type (b) for a dihedral group \(D_4=(\mathbf Z/2)^2\) otherwise.\par

(2.7) Lemma. Lemma 2.6 holds for a representation of dimension 2.\par

Such a representation is in fact either\par

(a) the realification of a one-dimensional complex representation,\par

(b) dihedral, hence the sum of a virtual representation \(r(\chi)\) and \([\varepsilon]+1\), for a character \(\varepsilon\) of order two.\par

(2.8) Lemma. Lemma 2.6 holds for a representation of \(G\) induced from a one-dimensional real representation of a subgroup of \(G\).\par

We proceed by induction on the index of the subgroup \(H\).\par

Let \(H'=G'\cap H\). If \(H\) has index \(\leq2\), apply Lemma 2.7. Otherwise, there is a subgroup \(J\) of \(G'\), belonging to the ascending central series, such that \(K=H'J/H'\) is nontrivial and abelian. Let \(l\) be a prime dividing the order of \(K\), and let \(K_l\) be the kernel of multiplication by \(l\). The group \(H/H'\) acts on \(K\) and \(K_l\).\par
}

\EditionPage{8}{306}{%
Since \(H/H'\) has order \(\leq2\), all its irreducible representations (over \(\mathbf Q\), or mod \(l\)) are one-dimensional, and the \(\mathbf F_l\)-vector space \(K_l\) contains a line stable under \(H/H'\). Let \(I'\) be its inverse image in \(H'J\). The group \(H\) normalizes \(I'\), so that \(I=I'H\) is a subgroup. We have \(I'=G'\cap I\), and \(H'\subset I'\subset G'\), with \(I'/H'\) cyclic of order \(l\).\par

Let \(\varepsilon:H\to\{\pm1\}\) be the character of the representation of \(H\) under consideration. The transitivity formula\par

\[
\operatorname{Ind}_H^G(\varepsilon)=\operatorname{Ind}_I^G\operatorname{Ind}_H^I(\varepsilon)
\]

and the induction hypothesis applied to \(I\) reduce us to showing that the representation \(\operatorname{Ind}_H^I(\varepsilon)\) of \(I\) satisfies Lemma 2.6. If \(l=2\), this representation has dimension two, and Lemma 2.7 applies. If \(l\neq2\), the nilpotence of \(G'\) ensures that \(I'\) normalizes \(\varepsilon|_{H'}\): otherwise, \(G'\) would have a subquotient that is a noncentral extension of a group of order \(l\) by a group of type \((2,\ldots,2)\), and such a group is not nilpotent.\par

Let us calculate the restriction of \(\operatorname{Ind}_H^I(\varepsilon)\) to \(I'\). It is \(\operatorname{Ind}_{H'}^{I'}(\varepsilon)\), obtained by inflation from \(\operatorname{Ind}_{H'/\ker(\varepsilon)}^{I'/\ker(\varepsilon)}(\varepsilon)\). The group \(I'/\ker(\varepsilon)\) is the product of a cyclic group of order \(l\) by \(H'/\ker(\varepsilon)\), of order \(\leq2\). Thus, as a complex representation, the representation of \(I'\) is a sum of distinct complex characters, and \(I/I'\) acts on the set of these characters either as the identity or by \(\chi\mapsto\chi^{-1}\) (according to its action on \(I'/H'\)). Grouping \(\chi\) and \(\chi^{-1}\) together, we find that the representation \(\operatorname{Ind}_H^I(\varepsilon)\) of \(I\) is a sum of real representations of dimension \(\leq2\), and Lemma 2.7 applies.\par

(2.9) Let us prove Lemma 2.6. We proceed by induction on the order of \(G\).\par

We may, and do, assume that\par

(1) \(V\) is an irreducible representation (over \(\mathbf R\));\par

(2) \(V\) is even absolutely irreducible (otherwise \(V\) is of type (a));\par

(3) \(V\) is faithful (otherwise the induction hypothesis is applied to a quotient of \(G\));\par

(4) \(V\) is not induced from a real representation of a proper subgroup of \(G\) (otherwise the induction hypothesis is applied to that subgroup, together with Lemma 2.8);\par

(5) \(G\) has no abelian subgroup of index \(\leq2\) (otherwise, by (2), \(\dim(V)\leq2\), and Lemma 2.7 applies).\par

Let \(B\) be an abelian normal subgroup of \(G\), contained in \(G'\), and maximal with these properties. If the complex character \(\chi\) of \(B\) occurs in the restriction of \(V\) to \(B\), then \(V\) is known to be induced from a real representation of the subgroup \(H\) of \(G\) that stabilizes \(\{\chi,\bar\chi\}\). By (1) and (4), \(H=G\), and only \(\chi\) and \(\bar\chi\) occur in \(V|_B\). Let \(G''\) be the stabilizer of \(\chi\). We have \(B\subset G'\cap G''\), and \(B\) is central in \(G''\).\par

Case 1. \(B\neq G'\cap G''\). By (3), \(B\) is central in the nilpotent group \(G'\cap G''\). The group \(G/(G'\cap G'')\) is of type \((\mathbf Z/2)^r\) (with \(r=0,1\), or 2); all its irreducible representations (over \(\mathbf Q\), or mod \(l\)) are one-dimensional. The argument at the beginning of the proof of Lemma 2.8 therefore applies again and shows that there is a subgroup \(B'\) of \(G'\cap G''\), normal in \(G\), with \(B'/B\) cyclic. This group \(B'\) is abelian, contradicting the maximality of \(B\).\par

Case 2. \(B=G'\cap G''\). Then \(B\) has index \(\leq2\) in \(G''\), which is therefore abelian. This contradicts (5).\par
}

\EditionPage{9}{307}{%
\EditionSection{3. Proof of 1.5}

(3.1) Return to the notation of 1.2 and 1.4. For \(V\) of dimension 0 and determinant 1, the number \(W(V)=\operatorname{sgn}(\varepsilon(V))\) is invariant under induction and multiplicative in \(V\). By 2.1 and local class field theory (namely, the formula \(\operatorname{inv}\operatorname{Tr}=\operatorname{inv}\)), \(\operatorname{cl}w^2(V)\) is likewise invariant under induction and additive in \(V\). Proposition 2.5 therefore reduces the proof of the theorem to the following two cases.\par

Case 1. \(V\) is the realification of a complex virtual representation \(W\) of dimension zero.\par

Let \(\chi:K^*\to\mathbf C^*\) be the determinant of \(W\). For every \(\psi\), \(\varepsilon(V)=\varepsilon(W,\psi)\cdot\varepsilon(\bar W,\psi)\); and\par

\[
\varepsilon(\bar W,\psi)=\chi(-1)\cdot\varepsilon(\bar W,\psi(-x))
=\chi(-1)\overline{\varepsilon(W,\psi)}\qquad\text{(1.2.2)},
\]

whence \(W(V)=\chi(-1)\).\par

In particular, \(W(V)=1\) when \(K\) has characteristic two, as it should. Now suppose that \(K\) has characteristic \(\neq2\). The class \(w^2(V)\) is the reduction modulo 2 of the first Chern class \(c^1(W)\). The class \(c^1(W)\in H^2(\operatorname{Gal}(\bar K/K),\mathbf Z)\) is, up to sign, the image of \(\chi\in H^1(\operatorname{Gal}(\bar K/K),\mathbf C^*)\) under the connecting homomorphism \(\partial\) in the long exact sequence arising from the exponential exact sequence\par

\[
0\longrightarrow\mathbf Z\xrightarrow{\,2\pi i\,}\mathbf C\xrightarrow{\exp}\mathbf C^*\longrightarrow0.
\]

The class \(w^2(V)\) is similarly obtained from \(\chi\) and the sequence\par

\[
0\longrightarrow\mathbf Z/2\longrightarrow\mathbf C^*\xrightarrow{z^2}\mathbf C^*\longrightarrow0.
\]

It is zero if and only if \(\chi\) has a square root, i.e., since \(-1\) is the unique element of order 2 in \(\widehat{K^*}=\operatorname{Gal}(\bar K/K)^{\mathrm{ab}}\), if and only if \(\chi(-1)=1\). This proves the theorem.\par

Case 2. \(V\) is obtained by inflation from a representation \(r(\chi)\) of a dihedral quotient of \(\operatorname{Gal}(\bar K/K)\).\par

For archimedean \(K\), necessarily \(V=0\). We may therefore assume, and do assume, that \(K\) is nonarchimedean.\par

Let \(G\) be a finite group that is an extension of \(\mathbf Z/2\) by \(\mathbf Z/n\). If \(u\in G\) does not belong to the subgroup \(\mathbf Z/n\), its image under the transfer \(G^{\mathrm{ab}}\to\mathbf Z/n\) is \(u^2\). Thus \(G\) is dihedral if and only if the transfer is trivial.\par

Recall also that, for an extension \(L\) of \(K\), the transfer \(\operatorname{Gal}(\bar K/K)^{\mathrm{ab}}\to\operatorname{Gal}(\bar K/L)^{\mathrm{ab}}\) corresponds to the inclusion \(K^*\to L^*\). The situation in Case 2 can therefore be described as follows: begin with a separable quadratic extension \(L/K\) and a character \(\chi:L^*\to\mathbf C^*\) trivial on \(K^*\). The representation \(\operatorname{Ind}\chi\) of \(\operatorname{Gal}(\bar K/K)\) is then realizable over \(\mathbf R\), is dihedral, and \(V=\operatorname{Ind}([\chi]-1)\).\par

(3.2) Theorem (Fröhlich and Queyrut [2]). For \(V\) as above, and for a nonzero element \(\Delta\) of \(L\) such that \(\operatorname{Tr}_{L/K}(\Delta)=0\) (\(\Delta\) is well defined up to multiplication by an element of \(K^*\)), one has\par

\[
W(V)=\chi(\Delta).
\]

Let \(\psi\) be a nontrivial additive character of \(K\), let \(\psi_L=\psi\circ\operatorname{Tr}_{L/K}\), and let \(dx\) be a Haar measure on \(L\). Then\par

\[
\begin{aligned}
\varepsilon(V,s)
&=\varepsilon(\operatorname{Ind}([\chi]-1),\psi,s)
 =\varepsilon([\chi]-1,\psi_L,s)\\
&=\varepsilon(\chi,\psi_L,dx,s)\cdot\varepsilon(1,\psi_L,dx,s)^{-1}.
\end{aligned}
\]
}

\EditionPage{10}{308}{%
Define the Fourier transform on the Schwartz-Bruhat functions on \(L\) by\par

\[
\widehat f(y)=\int f(x)\psi_L(xy)\,dx.
\]

The \(\varepsilon\)-factors are defined by Tate’s functional equation\par

\[
\tag{3.2.1}
\frac{\displaystyle\int \widehat f(x)\lVert x\rVert^{1-s}\chi^{-1}(x)\,d^*x}{L(\chi^{-1},1-s)}
=\varepsilon(\chi,\psi_L,dx,s)
\frac{\displaystyle\int f(x)\lVert x\rVert^s\chi(x)\,d^*x}{L(\chi,s)}
\]

(the same Haar measure \(d^*x\) on \(L^*\) is used on both sides). The two sides of this formula are generally defined by analytic continuation in \(s\). However, if \(f(0)=\widehat f(0)=0\), the integrals reduce to finite sums. Let \(x\mapsto\bar x\) be the nontrivial \(K\)-automorphism of \(L\). We have \(\chi^{-1}(x)=\chi(\bar x)\), so that \(L(\chi,s)=L(\chi^{-1},s)\). Setting \(s=\tfrac12\) in (3.2.1), we find that, when \(f(0)=\widehat f(0)=0\),\par

\[
\int \widehat f(x)\lVert x\rVert^{1/2}\chi^{-1}(x)\,d^*x
=\varepsilon\!\left(\chi,\psi_L,dx,\tfrac12\right)
\int f(x)\lVert x\rVert^{1/2}\chi(x)\,d^*x.
\]

We have\par

\[
\int f(x)\lVert x\rVert^{1/2}\chi(x)\,d^*x
=\int_{L^*/K^*}\lVert x\rVert^{1/2}\chi(x)(d^*x/d^*a)
\int_{K^*}f(ax)\,da,
\]

where \(da\) is a Haar measure on \(K\), and\par

\[
d^*a=\frac{da}{\lVert a\rVert_K}
\]

(note that \(\lVert a\rVert_L=\lVert a\rVert_K^2\)). Similarly,\par

\[
\begin{aligned}
\int \widehat f(x)\lVert x\rVert^{1/2}\chi^{-1}(x)\,d^*x
&=\int_{L^*/K^*}\lVert x\rVert^{1/2}\chi^{-1}(x)(d^*x/d^*a)
  \int_K\widehat f(ax)\,da\\
&=\int_{L^*/K^*}\lVert x\rVert^{1/2}\chi(x)(d^*x/d^*a)
  \int_K\widehat f(a\bar x)\,da.
\end{aligned}
\]

For a constant \(c\) depending on \(\psi_L,dx,da\), and \(\Delta\), one has\par

\[
\int_K\widehat f(a\bar x)\,da=c\int_K f(ax\Delta)\,da,
\]

so, putting \(c_1=c\cdot\lVert\Delta\rVert^{-1/2}\),\par

\[
\begin{aligned}
\int \widehat f(x)\lVert x\rVert^{1/2}\chi^{-1}(x)\,d^*x
&=\int_{L^*/K^*}\lVert x\rVert^{1/2}\chi(x)(d^*x/d^*a)c\cdot
  \int_K f(ax\Delta)\,da\\
&=c\cdot\lVert\Delta\rVert^{-1/2}\chi(\Delta)
  \int_{L^*/K^*}\lVert x\rVert^{1/2}\chi(x)(d^*x/d^*a)
  \int_K f(ax)\,da\\
&=c_1\chi(\Delta)\int f(x)\lVert x\rVert^{1/2}\chi(x)\,d^*x.
\end{aligned}
\]

For suitable choices of \(f\), the integrals can be made nonzero. Hence\par

\[
\varepsilon\!\left(\chi,\psi_L,dx,\tfrac12\right)=c_1\chi(\Delta)
\]

and\par

\[
W(V)=\varepsilon\!\left(V,\tfrac12\right)=\chi(\Delta).
\]

(3.3) Return to Case 2 in the proof of the theorem. If \(K\) has characteristic 2, then \(\Delta\in K^*\), and Theorem 3.2 shows that \(W(V)=1\), as it should.\par

Now suppose that \(K\) has characteristic \(\neq2\).\par
}

\EditionPage{11}{309}{%
The group \(O(2)\) is the semidirect product of \(\{e,s\}\) (with \(s\) a reflection) by \(SO(2)\). Let \(\widetilde O(2)\) be the semidirect product of \(\{e,s\}\) by the double cover of \(SO(2)\). For a representation \(\rho:G\to O(2)\) of the finite group \(G\), the obstruction to lifting \(\rho\) to \(\widetilde O(2)\) (the class of the central extension of \(G\) by \(\mathbf Z/2\) obtained by pulling back \(\widetilde O(2)\) under \(\rho\)) is of the form \(\alpha w_1^2+w_2\), with \(\alpha=0\) or 1. In fact, \(\alpha=0\), as is seen by testing on \(G=\{e,s\}\).\par

Put \(w^1=w^1(\operatorname{Ind}\chi)=w^1(\operatorname{Ind}1)\) and \(w^2=w^2(\operatorname{Ind}\chi)\). We have\par

\[
\begin{aligned}
w^*\operatorname{Ind}([\chi]-1)
&=(w^*\operatorname{Ind}(\chi))(w^*\operatorname{Ind}1)^{-1}\\
&=(1+w^1+w^2+\cdots)(1-w^1+(w^1)^2-\cdots)=1+w^2+\cdots.
\end{aligned}
\]

The class \(w^2\operatorname{Ind}([\chi]-1)\in H^2(\operatorname{Gal}(\bar K/K),\mathbf Z/2)\) is therefore zero if and only if the representation\par

\[
\operatorname{Ind}\chi:\operatorname{Gal}(\bar K/K)\longrightarrow O(2)
\]

lifts to \(\widetilde O(2)\). This means that \(\chi:L^*\to\mathbf C^*\) has a square root trivial on \(K^*\), or equivalently that \(\chi(\Delta)=1\), where \(\Delta\) is the unique element of order 2 in \(L^*/K^*\). This completes the proof of the theorem.\par

\EditionSection{4. Weil Groups and the Hasse Reciprocity Law}

(4.1) Let \(K\) be a local field and \(\bar K\) a separable closure of \(K\). For the definition of the (absolute) Weil group \(W(\bar K/K)\), see [1] 2.2. If \(K\) is nonarchimedean, it is a subgroup of \(\operatorname{Gal}(\bar K/K)\), an extension of \(\mathbf Z\) by inertia. Its cohomology with values in a discrete module \(M\) is defined by\par

\[
H^i(W(\bar K/K),M)=\varinjlim_U H^i(W(\bar K/K)/U,M^U)
\]

(\(U\) an open subgroup of the inertia group, normal in \(W(\bar K/K)\)). If \(K\) is archimedean, \(W(\bar K/K)\) is a Lie group, an extension of \(\operatorname{Gal}(\bar K/K)\) by \(\bar K^*\); its cohomology is that of the classifying space \(BW(\bar K/K)\).\par

Let \(n\) be an integer prime to the characteristic exponent of \(K\), and let \(\mu_n\) be the group of \(n\)-th roots of unity in \(\bar K\). It is a \(W(\bar K/K)\)-module.\par

(4.2) Construction. One constructs an isomorphism\par

\[
\operatorname{inv}:H^2(W(\bar K/K),\mu_n)\xrightarrow{\sim}(\mathbf Q/\mathbf Z)_n.
\]

(4.2.1) Lemma. If \(K\) is nonarchimedean and \(M\) is a discrete torsion \(\operatorname{Gal}(\bar K/K)\)-module, then\par

\[
H^i(\operatorname{Gal}(\bar K/K),M)\xrightarrow{\sim}H^i(W(\bar K/K),M).
\]

Let \(I\) be the inertia group. The map in 4.2.1 is the abutment of a morphism of spectral sequences\par

\[
E_2^{pq}=H^p(\widehat{\mathbf Z},H^q(I,M))\Rightarrow H^{p+q}(\operatorname{Gal}(\bar K/K),M),
\]

\[
E_2^{pq}=H^p(\mathbf Z,H^q(I,M))\Rightarrow H^{p+q}(W(\bar K/K),M),
\]

and one uses the fact that \(\mathbf Z\) and \(\widehat{\mathbf Z}\) have the same cohomology with values in a discrete torsion module (on which \(\widehat{\mathbf Z}\) acts continuously).\par

For nonarchimedean \(K\), therefore,\par

\[
H^2(\operatorname{Gal}(\bar K/K),\mu_n)\xrightarrow{\sim}H^2(W(\bar K/K),\mu_n).
\]
}

\EditionPage{12}{310}{%
One then uses the Kummer exact sequence\par

\[
0\longrightarrow\mu_n\longrightarrow\bar K^*\xrightarrow{n}\bar K^*\longrightarrow0
\]

and Theorem 90 to identify this \(H^2\) with the kernel of multiplication by \(n\) in\par

\[
H^2(\operatorname{Gal}(\bar K/K),\bar K^*)=\operatorname{Br}(K)=\mathbf Q/\mathbf Z.
\]

Given the state of the literature, this defines 4.2 up to sign. The sign will be discussed in 4.7. In the following section, in any case, we shall use only the case \(n=2\).\par

For archimedean \(K\), let \(\mathbf Z(1)=2\pi i\mathbf Z\subset\bar K\). We shall prove that \(H^3(W(\bar K/K),\mathbf Z(1))=0\) and construct an isomorphism\par

\[
\tag{4.2.2}\operatorname{inv}:H^2(W(\bar K/K),\mathbf Z(1))\xrightarrow{\sim}\mathbf Z.
\]

The long exact cohomology sequence associated with the short exact sequence\par

\[
0\longrightarrow\mathbf Z(1)\xrightarrow{n}\mathbf Z(1)
\xrightarrow{\exp(x/n)}\mu_n\longrightarrow0
\]

will give the required isomorphism\par

\[
\operatorname{inv}:H^2(W(\bar K/K),\mu_n)\xrightarrow{\sim}\mathbf Z/n
\xrightarrow[\,x\mapsto x/n\,]{\sim}(\mathbf Q/\mathbf Z)_n.
\]

If \(K\) is complex, then \(\bar K=K\), and \(H^2(W(\bar K/K),\mathbf Z(1))\) classifies extensions of \(K^*\) by \(\mathbf Z(1)\), i.e. homomorphisms from \(\pi_1(K^*)\) to \(\mathbf Z(1)\). This is the group \(\mathbf Z\), generated by the class of the extension\par

\[
\tag{4.2.3}0\longrightarrow\mathbf Z(1)\longrightarrow\bar K\xrightarrow{\exp}\bar K^*\longrightarrow0.
\]

The cohomology of \(W(\bar K/K)\) is that of infinite-dimensional complex projective space: the successive \(H^i\), with values in \(\mathbf Z(1)\), are\par

\[
\mathbf Z(1),0,\mathbf Z,0,\mathbf Z(1),\ldots.
\]

For real \(K\), use the spectral sequence\par

\[
E_2^{pq}=H^p(\mathbf Z/2,H^q(\bar K^*,\mathbf Z(1)))\Rightarrow H^{p+q}(W(\bar K/K),\mathbf Z(1)).
\]

\[
\begin{array}{c|cccccc}
H^q(\bar K^*,\mathbf Z(1)) & H^0H^q & H^1H^q & H^2H^q & H^3H^q & H^4H^q & \cdots\\ \hline
0            & 0         & 0           & 0           & 0           & 0           & \cdots\\
\mathbf Z    & \mathbf Z & 0           & \mathbf Z/2 & 0           & \mathbf Z/2 & \cdots\\
0            & 0         & 0           & 0           & 0           & 0           & \cdots\\
\mathbf Z(1) & 0         & \mathbf Z/2 & 0           & \mathbf Z/2 & 0           & \cdots
\end{array}
\]

We have \(E_2=E_3\), since \(E_2^{pq}=0\) for odd \(q\).\par

(4.2.4) Lemma. The homomorphism\par

\[
H^2(W(\bar K/K),\mathbf Z(1))\longrightarrow H^2(\bar K^*,\mathbf Z(1))
\]

is not surjective.\par

It is enough to see that extension 4.2.3 is not the restriction to \(\bar K^*\) of an extension \(E\) of \(W(\bar K/K)\) by \(\mathbf Z(1)\). The group \(W(\bar K/K)\) is generated by \(\bar K^*\) and \(F\), with relations \(F^2=-1\) and \(FzF^{-1}=\bar z\). If \(E\) existed, one could lift \(F\) to \(\widetilde F\) in \(E\). We would have \(\widetilde Fz\widetilde F^{-1}=\bar z\) (\(z\in\bar K\)) and \(\widetilde F^2\equiv\pi i\pmod{\mathbf Z(1)}\). It would follow that \(\widetilde F^2\) did not commute with \(\widetilde F\), which is absurd.\par

By the lemma, \(d_3:E_3^{0,2}\to E_3^{3,0}\) is nonzero, and the first \(H^i\) are\par

\[
H^0=0.
\]

\[
H^1=\mathbf Z/2\quad\text{(equal to }H^1(\operatorname{Gal}(\bar K/K),\mathbf Z(1))\text{)}.
\]
}

\EditionPage{13}{311}{%
\(H^2=\mathbf Z\), and \(H^2(W(\bar K/K),\mathbf Z(1))\hookrightarrow H^2(\bar K^*,\mathbf Z(1))\) as a subgroup of index two, the inclusion being identified with \(n\mapsto2n\).\par

\(H^3=0\).\par

(4.3) Remark. One can show that the cohomology is periodic with period 4. We shall not use this fact.\par

(4.4) Remark. Let \(K'\) be a finite extension of \(K\) in \(\bar K\), and let \(\alpha\in H^2(W(\bar K/K'),\mu_n)\). Then \(\operatorname{inv}\operatorname{Tr}\alpha=\operatorname{inv}\alpha\). This follows from the analogous statement for \(\operatorname{Br}(K)\) when \(K\) is nonarchimedean, and from the same formula for \(\alpha\in H^2(W(\bar K/K'),\mathbf Z(1))\) when \(K\) is archimedean. The latter follows from the formulas \(\operatorname{Tr}\operatorname{res}\beta=2\beta\) and \(\operatorname{inv}\operatorname{res}\beta=2\operatorname{inv}\beta\) for \(\beta\in H^2(W(\bar K/K),\mathbf Z(1))\), with \(K\) real and \(K'\) complex.\par

(4.5) Example. Let \(\chi:K^*\to\mathbf C^*\) be a quasicharacter. The class \(\partial\chi\) of the extension of \(K^*\) by \(\mathbf Z\) obtained by pulling back along \(\chi\) the extension\par

\[
\tag{4.5.1}0\longrightarrow\mathbf Z\xrightarrow{2\pi i}\mathbf C\xrightarrow{\exp}\mathbf C^*\longrightarrow0
\]

is an element of \(H^2(K^*,\mathbf Z)\). If \(\alpha\) is a root of unity in \(K\) whose order divides \(n\), its image under \(\mathbf Z\to\mu_n:m\mapsto\alpha^m\) gives\par

\[
(\chi,\alpha)\in H^2(W(\bar K/K),\mu_n).
\]

(4.6) Proposition. We have\par

\[
\exp\bigl(2\pi i\,\operatorname{inv}(\chi,\alpha)\bigr)=\chi(\alpha).
\]

The archimedean case is verified by a direct calculation. When \(K\) is nonarchimedean, neither side changes if \(\chi\) is multiplied by an unramified character. We may therefore suppose that \(\chi\) has finite order. It then extends to\par

\[
\chi:\operatorname{Gal}(\bar K/K)^{\mathrm{ab}}\longrightarrow\mathbf C^*,
\]

which defines \([\chi]\in H^1(\operatorname{Gal}(\bar K/K),\mathbf C^*)\). Its image \(\partial\chi\in H^2(\operatorname{Gal}(\bar K/K),\mathbf Z)\), under the connecting homomorphism \(\partial\) associated with (4.5.1), is the class of the extension of \(\operatorname{Gal}(\bar K/K)\) by \(\mathbf Z\) obtained by pulling back (4.5.1) along \(\chi\). Its image under \(\mathbf Z\to\mu_n\), \(m\mapsto\alpha^m\), is the cup product \(\partial\chi\cup\alpha\in H^2(\operatorname{Gal}(\bar K/K),\mu_n)\). Its restriction to \(W(\bar K/K)\) is \((\chi,\alpha)\), and \(\operatorname{inv}(\chi,\alpha)\) is the invariant of the cup product \(\partial\chi\cup\alpha\in H^2(\operatorname{Gal}(\bar K/K),\bar K^*)\). In this form, the formula\par

\[
\exp\bigl(2\pi i\,\operatorname{inv}(\partial\chi\cup\alpha)\bigr)=\chi(\alpha)
\]

holds for every \(\alpha\in K^*\): this is the definition of the reciprocity isomorphism.\par

(4.7) Signs. Since we have already specified our normalization (= choice of sign) of the reciprocity isomorphism, the validity of 4.6 fixes our normalization of \(\operatorname{inv}\) for nonarchimedean \(K\).\par

(4.8) Let \(K\) be a global field. For a finite Galois extension \(L\) of \(K\), Weil defined the Weil group \(W_{L/K}\), an extension of \(\operatorname{Gal}(L/K)\) by the idèle class group \(C(L)\) of \(L\). It is an inverse limit of Lie groups; its cohomology is by definition the direct limit of the cohomologies of their classifying spaces. Let \(\bar K\) be a separable closure of \(K\), and let \(M\) be a discrete \(\operatorname{Gal}(\bar K/K)\)-module. We set\par

\[
H^*(W(\bar K/K),M)=\varinjlim_{L\subset\bar K}H^*(W_{L/K},M^{\operatorname{Gal}(\bar K/L)}).
\]
}

\EditionPage{14}{312}{%
Let \(n\) be prime to the characteristic exponent of \(K\), and let \(\alpha\in H^2(W(\bar K/K),\mu_n)\). The restriction of \(\alpha\) to \(H^2(W(\bar K_v/K_v),\mu_n)\) is zero for almost all places \(v\) of \(K\) (because it factors through an \(H^2(\mathbf Z,\mu_n)=0\)).\par

The maps \(\operatorname{inv}\) of 4.2 therefore give a map\par

\[
\operatorname{inv}:H^2(W(\bar K/K),\mu_n)\longrightarrow\bigoplus_v(\mathbf Q/\mathbf Z)_n.
\]

(4.9) Theorem. The sequence\par

\[
0\longrightarrow H^2(W(\bar K/K),\mu_n)
\xrightarrow{\operatorname{inv}}\bigoplus_v(\mathbf Q/\mathbf Z)_n
\xrightarrow{\Sigma}(\mathbf Q/\mathbf Z)_n\longrightarrow0
\]

is exact.\par

Let \(\chi:C(K)\to\mathbf C^*\) be a quasicharacter. The class \(\partial\chi\) of the extension of \(C(K)\) by \(\mathbf Z\) obtained by pulling back along \(\chi\) the extension\par

\[
0\longrightarrow\mathbf Z\xrightarrow{2\pi i}\mathbf C\xrightarrow{\exp}\mathbf C^*\longrightarrow0
\]

is an element of \(H^2(C(K),\mathbf Z)\). If \(\alpha\) is a root of unity in \(K\) whose order divides \(n\), its image under \(\mathbf Z\to\mu_n:m\mapsto\alpha^m\) gives\par

\[
(\chi,\alpha)\in H^2(C(K),\mu_n).
\]

(4.9.1) Lemma. We have\par

\[
\sum_v\operatorname{inv}_v(\chi,\alpha)=0.
\]

By 4.6, this is a reformulation of \(\prod_v\chi_v(\alpha_v)=1\).\par

(4.9.2) Lemma. There exists \(c\in H^2(W(\bar K/K),\mu_n)\) having prescribed local invariants at the infinite places (or, indeed, at any finite set of places) and satisfying\par

\[
\sum_v\operatorname{inv}_v(c)=0.
\]

If \(\mu_n\subset K\), one may take \(c\) of the form \((\chi,\alpha)\). Otherwise, take a finite extension \(K'/K\) such that \(K'\) contains \(\mu_n\), and take \(c\) of the form \(\operatorname{Tr}_{K'/K}(\chi,\alpha)\); the reciprocity law follows from 4.4 and 4.9.1.\par

(4.9.3) Lemma. Let \(c\in H^2(W(\bar K/K),\mu_n)\). If \(\operatorname{inv}_v(c)=0\) for complex \(v\), and has order dividing two for real \(v\), then \(c\) comes from \(\bar c\in H^2(\operatorname{Gal}(\bar K/K),\mu_n)\).\par

When \(K\) has nonzero characteristic, one passes from a Weil group to a Galois group by replacing a quotient \(\mathbf Z\) by \(\widehat{\mathbf Z}\), and it is enough to argue as in 4.2.1.\par

Now suppose that \(K\) is a number field. There exists a finite totally imaginary Galois extension \(L\) of \(K\), containing \(\mu_n\) if desired, such that \(c\) is defined by an extension of \(W_{L/K}\) by \(\mu_n\). There is even a modulus \(\mathfrak m\) such that \(c\) comes from an extension \(E'\), by \(\mu_n\), of the quotient \({}_{\mathfrak m}W_{L/K}\) of \(W_{L/K}\), which is an extension of \(\operatorname{Gal}(L/K)\) by \(C_{\mathfrak m}(L)\). The identity component of \(C_{\mathfrak m}(L)\) is the quotient \(L_{\mathbf R}^*/U_{\mathfrak m}\) of \((L\otimes\mathbf R)^*\) by the group of units congruent to 1 modulo \(\mathfrak m\). The hypothesis ensures that the extension induced by \(E'\) of \(L_{\mathbf R}^*\) by \(\mu_n\) is trivial. The extension of \(L_{\mathbf R}^*/U_{\mathfrak m}\) by \(\mu_n\) therefore comes from \(\phi:U_{\mathfrak m}\to\mu_n\). By a theorem of Chevalley (deux théorèmes d’arithmétique, J. Math. Soc. Jap. 3 (1951), pp. 36--44), \(\operatorname{Ker}(\phi)\) is a congruence subgroup: for a suitable modulus \(\mathfrak m'\ge\mathfrak m\), \(c\) is defined by an extension \(E\) of \({}_{\mathfrak m'}W_{L/K}\) by \(\mu_n\), trivial on the identity component of \({}_{\mathfrak m'}W_{L/K}\). The extension \(E\) therefore comes from the Galois group \(\pi_0({}_{\mathfrak m'}W_{L/K})\); this proves 4.9.3.\par
}

\EditionPage{15}{313}{%
Set\par

\[
(\mathbf Q/\mathbf Z)_{n,v}=
\begin{cases}
\text{the subgroup of }\mathbf Q/\mathbf Z\text{ of order }n,&v\text{ nonarchimedean},\\
\text{the subgroup of order }(n,2),&v\text{ real},\\
\text{the subgroup of order }1,&v\text{ complex},
\end{cases}
\]

and\par

\[
A=\left(\bigoplus_v(\mathbf Q/\mathbf Z)_n\right)\Big/\left(\bigoplus_v(\mathbf Q/\mathbf Z)_{n,v}\right).
\]

The proof of 4.9 is completed by considering the following diagram, in which exactness of the last row is the Hasse reciprocity law:\par

\[
\begin{array}{ccccccc}
&&0&&&&\\
&&\uparrow&&&&\\
A&=&A&&&&\\
\uparrow&&\uparrow&&&&\\
H^2(W(\bar K/K),\mu_n)&\xrightarrow{\operatorname{inv}}&\displaystyle\bigoplus_v(\mathbf Q/\mathbf Z)_n&\xrightarrow{\Sigma}&(\mathbf Q/\mathbf Z)_n&\longrightarrow&0\\
\uparrow&&\uparrow&&\Vert&&\\
0\longrightarrow H^2(\operatorname{Gal}(\bar K/K),\mu_n)&\xrightarrow{\operatorname{inv}}&\displaystyle\bigoplus_v(\mathbf Q/\mathbf Z)_{n,v}&\xrightarrow{\Sigma}&(\mathbf Q/\mathbf Z)_n&\longrightarrow&0\\
&&\uparrow&&&&\\
&&0&&&&
\end{array}
\]

\EditionSection{5. Orthogonal representations of Weil groups}

(5.1) Rather than considering real representations of Galois groups, it is more general and more natural to consider complex orthogonal representations of Weil groups.\par

Let \(K\) be a local field, let \(\bar K\) be a separable closure of \(K\), let \(V\) be a complex vector space equipped with a nondegenerate quadratic form \(Q\), and let \(\rho:W(\bar K/K)\to O(V,Q)\) be a continuous orthogonal representation of \(W(\bar K/K)\). The Stiefel--Whitney classes \(w^i\in H^i(W(\bar K/K),\mathbf Z/2)\) are still defined. Since the space of nondegenerate quadratic forms on \(V\) invariant under \(W(\bar K/K)\) is connected, they depend only on the linear representation underlying \(\rho\). The total Stiefel--Whitney class has the following additivity properties:\par

(5.1.1) If \(V\) is the orthogonal direct sum of subrepresentations \(V'\) and \(V''\), then\par

\[
w^*(V)=w^*(V')\cdot w^*(V'').
\]

(5.1.2) If \(W\) is an invariant totally isotropic subspace of \(V\), then\par

\[
w^*(V)=w^*(W\oplus V/W^\perp)\cdot w^*(W^\perp/W).
\]

The quadratic form \(Q\) on \(V\) induces a duality between \(W\) and \(V/W^\perp\)---hence a quadratic form on \(W\oplus V/W^\perp\)---whereas \(Q|W^\perp\) has kernel \(W\)---hence a quadratic form on \(W^\perp/W\). This gives meaning to 5.1.2. To prove the formula, it is enough to prove the analogous statement for a complex orthogonal vector bundle \(\mathscr V\) on a paracompact topological space and for \(\mathscr W\subset\mathscr V\)\par
}

\EditionPage{16}{314}{%
totally isotropic. Here \(\mathscr V\) is isomorphic to \((\mathscr W\oplus\mathscr V/\mathscr W^\perp)\oplus(\mathscr W^\perp/\mathscr W)\), and one applies the analogue of 5.1.1.\par

Let \(RO(W(\bar K/K),\mathbf C)\) be the Grothendieck group of semisimple complex orthogonal representations of \(W(\bar K/K)\), with respect to orthogonal direct sum. It is a subgroup of \(R(W(\bar K/K),\mathbf C)\), the Grothendieck group of the category of semisimple complex representations of \(W(\bar K/K)\). Property (5.1.1) makes it possible to define the total Stiefel--Whitney class of an element of \(RO(W(\bar K/K),\mathbf C)\) (a virtual orthogonal representation).\par

If \(V\) is an orthogonal representation of \(W(\bar K/K)\), its semisimplification \(V'\) is again orthogonal, and (5.1.2) gives \(w^*(V)=w^*(V')\). This justifies the restriction above to semisimple representations.\par

For \(x\in H^2(W(\bar K/K),\mathbf Z/2)\), the class \(\operatorname{cl}(x)\) is 0 when \(K\) has characteristic 2, and otherwise is the invariant of its image in \(H^2(W(\bar K/K),\mu_2)\) (with \(\mathbf Z/2\) and \(\mu_2\) canonically isomorphic). When \(K\) is not complex, one has\par

\[
H^2(\operatorname{Gal}(\bar K/K),\mathbf Z/2)\xrightarrow{\sim}H^2(W(\bar K/K),\mathbf Z/2),
\]

and this definition is compatible with 1.4. For \(K\) of characteristic different from 2, archimedean or not, \(\operatorname{cl}\) is an isomorphism\par

\[
\operatorname{cl}:H^2(W(\bar K/K),\mathbf Z/2)\xrightarrow{\sim}(\mathbf Q/\mathbf Z)_2.
\]

For a complex virtual representation \(V\) of \(W(\bar K/K)\), of dimension 0 and determinant 1 and isomorphic to its dual, define, as in 1.2, a number \(W(V)=\pm1\). Again one has:\par

(5.2) Proposition. Let \(V\) be a virtual orthogonal representation of dimension 0 and determinant 1. Then\par

\[
W(V)=\exp\bigl(2\pi i\,\operatorname{cl}(w^2(V))\bigr).
\]

For \(V\) of the form \(W\oplus W^*\) (where \(W\) is a virtual representation of dimension 0 and determinant \(\chi\)), one proceeds as in 3.1, Case 1, using the formulas\par

\[
w^2(V)=c^1(W)\pmod2
\]

and\par

\[
\varepsilon(W,\psi)\cdot\varepsilon(W^*(1),\psi(-x))=1.
\]

This suffices to prove 5.2 when \(K\) is complex. When \(K\) is real, it remains to treat the dihedral case of a representation \(r(\chi)\), where \(\chi\) is a character of \(\bar K^*\) trivial on \(K^*\). Theorem 3.2 remains valid, with the same proof [this time, for \(f\) and \(\hat f\), smooth and rapidly decreasing and having a zero of high multiplicity at 0, the integrals converge in a wide strip], and one repeats the proof of 3.3.\par

When \(K\) is nonarchimedean, one verifies using ([1] 4.10) that every semisimple orthogonal representation \(V\) of \(W(\bar K/K)\) is the sum of a representation of the form \(W\oplus W^*\) and a representation factoring through a finite quotient of \(\operatorname{Gal}(\bar K/K)\). For virtual \(V\) of dimension 0, one obtains a similar decomposition with \(W\) of dimension 0, and the proposition therefore follows from Theorem 1.5.\par

(5.3) Corollary. Let \(K\) be a global field, let \(L\) be a finite extension of \(K\), and let \(V\) be an orthogonal representation of the Weil group \(W_{L/K}\). The constant \(W(V)\) in the functional equation of the corresponding \(L\)-function, defined as in (1.1.2), is \(+1\).\par

The proof proceeds as in 1.6, now using 5.2 and the reciprocity law 4.9 (for \(n=2\)).\par

(5.4) \(\ell\)-adic variant. Let \(K\) be a nonarchimedean local field, and let \(V\) be a complex virtual representation of \(W(\bar K/K)\), of dimension 0, with trivial determinant, and\par
}

\EditionPage{17}{315}{%
isomorphic to its dual. The definition of \(\varepsilon(V)=\varepsilon(V,0)\) is purely algebraic, i.e. independent of the topology of \(\mathbf C\). If the conductor of \(V\) is not a square, the same is not true of \(W(V)=\varepsilon(V)\cdot f^{-1/2}\). However, if \(V\) is orthogonal, Serre proved [3] that the conductor exponent is even; in this case \(W(V)\) therefore has a purely algebraic meaning. The same holds for \(w^2(V)\). By transport of structure, it follows that the formula (to which a meaning has just been given)\par

\[
\tag{5.4.1}W(V)=(-1)^{2\operatorname{cl}(w^2(V))}
\]

holds for a virtual orthogonal representation \(V\) of \(W(\bar K/K)\), of dimension 0 and trivial determinant, over any algebraically closed field \(k\) of characteristic 0. The hypothesis “algebraically closed” is unnecessary, as is seen by replacing \(k\) with its algebraic closure.\par

(5.5) Let \((V,\rho)\) be a continuous \(\ell\)-adic representation of \(W(\bar K/K)\), and fix \(\psi\) and \(dx\). Following [1] 8.4.2, one associates with it a “representation” \((\rho',N)\) of \({}'W(\bar K/K)\), consisting of\par

(a) a representation \(\rho'\) of \(W(\bar K/K)\) on \(V\), continuous for the discrete topology on \(V\);\par

(b) a nilpotent endomorphism \(N\) of \(V\), such that\par

\[
\rho'(w)N\rho'(w)^{-1}=q^{v(w)}N.
\]

If \(\rho\) preserves a quadratic form \(Q\), then so does \(\rho'\), and \(N\in\operatorname{Lie}(O(Q))\).\par

The constant \(\varepsilon(V,\psi,dx,t)\) (where \(t\) is morally \(q^{-s}\)), defined by [1] 8.12, is the product of \(\varepsilon(\rho',\psi,dx,t)\) and\par

\[
\det\bigl(-Ft,\,V^{\rho'(I)}/\operatorname{Ker}(N)^{\rho'(I)}\bigr).
\]

This monomial depends only on \(V^{\rho'(I)}\), on which \(F\) and \(N\) act (with \(FNF^{-1}=q^{-1}N\)). The lemma below shows that, if \(V\) is orthogonal, it has even degree and takes the value 1 at \(t=q^{-1/2}\).\par

For a virtual \(\ell\)-adic orthogonal representation \(V\) of degree 0 and determinant one, the conductor of \(V\) therefore still has even exponent, and \(W(V)\) may be defined as before. Moreover, \(W(V)=W(V^{\mathrm{ss}})\), where \(V^{\mathrm{ss}}\) is the semisimplification of \(V\). This result may be used to calculate the right-hand side of (1.2.1) for an \(\ell\)-adic orthogonal representation \(V\) of the Weil group of a global field; the method of [2] § 9 calculates the left-hand side, and one finds that (1.2.1) remains valid in this case.\par

(5.6) Lemma. Let \(k\) be a field, let \(q\in k^*\), let \(V\) be a vector space over \(k\), let \(Q\) be a nondegenerate quadratic form on \(V\), and let \(F\in O(V,Q)\) and \(N\in\operatorname{Lie}(O(Q))\). Suppose that \(FNF^{-1}=q^{-1}N\). Then \(\dim(V/\operatorname{Ker}N)\) is an even integer \(2m\), and\par

\[
\det(-Ft,V/\operatorname{Ker}N)=q^m t^{2m}.
\]

Let \(B\) be the bilinear form associated with \(Q\). By hypothesis, the form \(\langle x,y\rangle=B(Nx,y)\) is alternating, and \(\langle Fx,Fy\rangle=q\langle x,y\rangle\). The kernel of \(\langle x,y\rangle\) is \(\operatorname{Ker}N\), and \(F\) induces a symplectic similitude of multiplier \(q\) on \(V/\operatorname{Ker}N\). The lemma follows.\par
}

\EditionPage{18}{316}{%
\EditionHeading{Bibliography}

\noindent\hangindent=1.5em\hangafter=1 1. Deligne, P.: Les constantes des équations fonctionnelles des fonctions \(L\). In: Modular functions of one variable II, Lecture Notes in Mathematics 349, pp. 501--597, Berlin-Heidelberg-New York: Springer 1973\par

\noindent\hangindent=1.5em\hangafter=1 2. Fröhlich, A., Queyrut, J.: On the functional equation of the Artin \(L\)-function for characters of real representations, Inventiones math. 20, 125--138 (1973)\par

\noindent\hangindent=1.5em\hangafter=1 3. Serre, J-P.: Conducteurs d’Artin des caractères réels. Inventiones math. 14, 173--183 (1971)\par

\medskip\noindent\itshape Received 20 September 1975\par\normalfont

\medskip\noindent Pierre Deligne\\I.H.E.S.\\35, Route de Chartres\\F-91440 Bures-sur-Yvette/France\par
}

\end{document}
